\chapter{Inecuaciones}

\begin{theorybox}{Trazabilidad del capitulo}
Este capitulo desarrolla las secciones \texttt{C04.S01-C04.S07} de la matriz de cobertura a
partir de \texttt{sources/Relacion Tema 4. Inecuaciones..pdf}. El acento didactico esta en
traducir correctamente la solucion a intervalos y en justificar los signos.
\end{theorybox}

\section{[C04.S01] Inecuaciones lineales}

\subsection*{Objetivos de aprendizaje}
\begin{itemize}
  \item Resolver inecuaciones lineales de una variable.
  \item Expresar la solucion con desigualdades e intervalos.
\end{itemize}

\subsection*{Prerrequisitos}
\begin{itemize}
  \item Ecuaciones lineales y orden en \(\mathbb{R}\).
\end{itemize}

\begin{theorybox}{Idea minima}
Una inecuacion lineal se resuelve igual que una ecuacion lineal, con una precaucion importante:
si multiplicamos o dividimos por un numero negativo, el signo de la desigualdad cambia.
\end{theorybox}

\begin{methodbox}{Como reconocer y resolver}
\begin{enumerate}
  \item Reune los terminos con \(x\) en un miembro y los numeros en el otro.
  \item Simplifica.
  \item Si divides por un numero negativo, invierte el signo.
  \item Expresa el resultado en desigualdad o intervalo.
\end{enumerate}
\end{methodbox}

\begin{solvedexample}{[EX-C04.S01-01] Inecuacion lineal directa}
Resuelve
\[
3x-5\leq 7.
\]

\textbf{Desarrollo.}
\[
3x\leq 12 \Longrightarrow x\leq 4.
\]

\textbf{Respuesta final.}
\[
x\leq 4 \qquad \text{o bien} \qquad (-\infty,4].
\]
\end{solvedexample}

\begin{commonerror}{Olvidar cambiar el signo}
Si al despejar aparece una division por un numero negativo, el sentido de la desigualdad debe
invertirse.
\end{commonerror}

\begin{guidedexercise}{[GX-C04.S01-01] Un paso de suma}
Resuelve
\[
2x+3>7.
\]
\end{guidedexercise}

\begin{shortanswer}{[GX-C04.S01-01]}
\[
x>2.
\]
\end{shortanswer}

\begin{fullsolution}{[GX-C04.S01-01]}
\[
2x>4 \Longrightarrow x>2.
\]
\end{fullsolution}

\begin{guidedexercise}{[GX-C04.S01-02] Division por negativo}
Resuelve
\[
-3x+6\geq 0.
\]
\end{guidedexercise}

\begin{shortanswer}{[GX-C04.S01-02]}
\[
x\leq 2.
\]
\end{shortanswer}

\begin{fullsolution}{[GX-C04.S01-02]}
\[
-3x\geq -6.
\]
Al dividir por \(-3\), el signo cambia:
\[
x\leq 2.
\]
\end{fullsolution}

\subsection*{Practica autonoma graduada}
\begin{exercise}{[PX-C04.S01-01] Resuelve \(x-4<0\).}
\end{exercise}
\begin{exercise}{[PX-C04.S01-02] Resuelve \(5x+1\geq 11\).}
\end{exercise}
\begin{exercise}{[PX-C04.S01-03] Resuelve \(7-2x>1\).}
\end{exercise}
\begin{exercise}{[PX-C04.S01-04] Resuelve \(4x-3\leq 5\).}
\end{exercise}
\begin{exercise}{[PX-C04.S01-05] Resuelve \(2-3x<11\).}
\end{exercise}
\begin{exercise}{[PX-C04.S01-06] Resuelve \(-x+8\geq 10\).}
\end{exercise}

\begin{shortanswer}{C04.S01 practica autonoma}
\begin{enumerate}[label=\textbf{PX-C04.S01-\arabic*:}, leftmargin=*]
  \item \(x<4\)
  \item \(x\geq 2\)
  \item \(x<3\)
  \item \(x\leq 2\)
  \item \(x>-3\)
  \item \(x\leq -2\)
\end{enumerate}
\end{shortanswer}

\begin{fullsolution}{C04.S01 practica autonoma}
\begin{enumerate}[label=\textbf{PX-C04.S01-\arabic*:}, leftmargin=*]
  \item \(x<4\).
  \item \(5x\geq 10\), luego \(x\geq 2\).
  \item \(-2x>-6\), al dividir por \(-2\), \(x<3\).
  \item \(4x\leq 8\), luego \(x\leq 2\).
  \item \(-3x<9\), al dividir por \(-3\), \(x>-3\).
  \item \(-x\geq 2\), luego \(x\leq -2\).
\end{enumerate}
\end{fullsolution}

\section{[C04.S02] Inecuaciones de segundo grado}

\subsection*{Objetivos de aprendizaje}
\begin{itemize}
  \item Resolver inecuaciones cuadraticas.
  \item Usar las raices para estudiar el signo del polinomio.
\end{itemize}

\subsection*{Prerrequisitos}
\begin{itemize}
  \item Factorizacion de polinomios de segundo grado.
\end{itemize}

\begin{theorybox}{Idea minima}
Una parabola cambia de signo en sus raices. Si el coeficiente principal es positivo, el
polinomio es positivo fuera de las raices y negativo entre ellas; si es negativo, ocurre al
reves.
\end{theorybox}

\begin{methodbox}{Como reconocer y resolver}
\begin{enumerate}
  \item Factoriza o halla las raices.
  \item Situa las raices en la recta real.
  \item Decide en que intervalos el polinomio cumple la desigualdad.
\end{enumerate}
\end{methodbox}

\begin{solvedexample}{[EX-C04.S02-01] Cuadratica factorizada}
Resuelve
\[
x^2-5x+6>0.
\]

\textbf{Desarrollo.} Factorizamos:
\[
x^2-5x+6=(x-2)(x-3).
\]
Las raices son \(2\) y \(3\). Como el coeficiente de \(x^2\) es positivo, el producto es
positivo fuera del intervalo entre raices.

\textbf{Respuesta final.}
\[
x\in (-\infty,2)\cup (3,+\infty).
\]
\end{solvedexample}

\begin{commonerror}{Confundir \(>\) con \(\geq\)}
Si la desigualdad es estricta, las raices no se incluyen en la solucion.
\end{commonerror}

\begin{guidedexercise}{[GX-C04.S02-01] Entre dos raices}
Resuelve
\[
x^2-9\leq 0.
\]
\end{guidedexercise}

\begin{shortanswer}{[GX-C04.S02-01]}
\[
[-3,3].
\]
\end{shortanswer}

\begin{fullsolution}{[GX-C04.S02-01]}
\[
x^2-9=(x-3)(x+3).
\]
El producto es no positivo entre las raices:
\[
x\in [-3,3].
\]
\end{fullsolution}

\begin{guidedexercise}{[GX-C04.S02-02] Factor comun en el signo}
Resuelve
\[
x^2-4x<0.
\]
\end{guidedexercise}

\begin{shortanswer}{[GX-C04.S02-02]}
\[
(0,4).
\]
\end{shortanswer}

\begin{fullsolution}{[GX-C04.S02-02]}
\[
x^2-4x=x(x-4).
\]
El producto es negativo entre \(0\) y \(4\):
\[
x\in (0,4).
\]
\end{fullsolution}

\subsection*{Practica autonoma graduada}
\begin{exercise}{[PX-C04.S02-01] Resuelve \(x^2-1\geq 0\).}
\end{exercise}
\begin{exercise}{[PX-C04.S02-02] Resuelve \(x^2-6x+8<0\).}
\end{exercise}
\begin{exercise}{[PX-C04.S02-03] Resuelve \(x^2+4x+3\leq 0\).}
\end{exercise}
\begin{exercise}{[PX-C04.S02-04] Resuelve \(x^2-9>0\).}
\end{exercise}
\begin{exercise}{[PX-C04.S02-05] Resuelve \(x^2-2x-3\geq 0\).}
\end{exercise}
\begin{exercise}{[PX-C04.S02-06] Resuelve \(-x^2+4x-4\leq 0\).}
\end{exercise}

\begin{shortanswer}{C04.S02 practica autonoma}
\begin{enumerate}[label=\textbf{PX-C04.S02-\arabic*:}, leftmargin=*]
  \item \((-\infty,-1]\cup [1,+\infty)\)
  \item \((2,4)\)
  \item \([-3,-1]\)
  \item \((-\infty,-3)\cup (3,+\infty)\)
  \item \((-\infty,-1]\cup [3,+\infty)\)
  \item \(\mathbb{R}\)
\end{enumerate}
\end{shortanswer}

\begin{fullsolution}{C04.S02 practica autonoma}
\begin{enumerate}[label=\textbf{PX-C04.S02-\arabic*:}, leftmargin=*]
  \item \((x-1)(x+1)\geq 0\) fuera de las raices.
  \item \((x-2)(x-4)<0\) entre \(2\) y \(4\).
  \item \((x+1)(x+3)\leq 0\) entre las raices e incluyendolas.
  \item \(x^2-9>0\) fuera de \(\pm 3\).
  \item \((x-3)(x+1)\geq 0\) fuera del intervalo interior.
  \item \(-x^2+4x-4=-(x-2)^2\leq 0\) para todo real.
\end{enumerate}
\end{fullsolution}

\section{[C04.S03] Inecuaciones polinomicas}

\subsection*{Objetivos de aprendizaje}
\begin{itemize}
  \item Resolver inecuaciones polinomicas de grado mayor que dos.
  \item Usar tablas de signos apoyadas en los factores.
\end{itemize}

\subsection*{Prerrequisitos}
\begin{itemize}
  \item Factorizacion y estudio de signo en productos.
\end{itemize}

\begin{theorybox}{Idea minima}
En una inecuacion polinomica factorizada, los puntos clave son las raices de cada factor. Entre
dos puntos consecutivos, el signo del producto es constante.
\end{theorybox}

\begin{methodbox}{Como reconocer y resolver}
\begin{enumerate}
  \item Factoriza el polinomio.
  \item Ordena las raices reales.
  \item Estudia el signo en cada intervalo.
  \item Incluye o excluye las raices segun el signo de desigualdad.
\end{enumerate}
\end{methodbox}

\begin{solvedexample}{[EX-C04.S03-01] Tres factores lineales}
Resuelve
\[
(x-1)(x+2)(x-4)\geq 0.
\]

\textbf{Desarrollo.} Las raices son \(-2\), \(1\) y \(4\). El estudio de signos da:
\[
(-\infty,-2): -, \qquad (-2,1): +, \qquad (1,4): -, \qquad (4,+\infty): +.
\]
Como queremos \(\geq 0\), tomamos los intervalos positivos y las raices.

\textbf{Respuesta final.}
\[
[-2,1]\cup [4,+\infty).
\]
\end{solvedexample}

\begin{commonerror}{No tener en cuenta la multiplicidad}
Si una raiz tiene multiplicidad par, el signo no cambia al pasar por ella. Ese detalle altera la
tabla de signos.
\end{commonerror}

\begin{guidedexercise}{[GX-C04.S03-01] Dos factores}
Resuelve
\[
x(x-3)\leq 0.
\]
\end{guidedexercise}

\begin{shortanswer}{[GX-C04.S03-01]}
\[
[0,3].
\]
\end{shortanswer}

\begin{fullsolution}{[GX-C04.S03-01]}
Las raices son \(0\) y \(3\). El producto es no positivo entre ellas:
\[
[0,3].
\]
\end{fullsolution}

\begin{guidedexercise}{[GX-C04.S03-02] Raiz doble}
Resuelve
\[
(x+1)^2(x-2)>0.
\]
\end{guidedexercise}

\begin{shortanswer}{[GX-C04.S03-02]}
\[
(2,+\infty).
\]
\end{shortanswer}

\begin{fullsolution}{[GX-C04.S03-02]}
El factor \((x+1)^2\) es siempre no negativo y solo vale \(0\) en \(x=-1\). El signo del
producto lo decide \(x-2\), por lo que solo es positivo si \(x>2\).
\end{fullsolution}

\subsection*{Practica autonoma graduada}
\begin{exercise}{[PX-C04.S03-01] Resuelve \(x(x-1)(x+1)<0\).}
\end{exercise}
\begin{exercise}{[PX-C04.S03-02] Resuelve \((x-2)^2(x+3)\geq 0\).}
\end{exercise}
\begin{exercise}{[PX-C04.S03-03] Resuelve \(x(x-4)^2\leq 0\).}
\end{exercise}
\begin{exercise}{[PX-C04.S03-04] Resuelve \((x+2)(x-5)>0\).}
\end{exercise}
\begin{exercise}{[PX-C04.S03-05] Resuelve \(x^2(x-1)\leq 0\).}
\end{exercise}
\begin{exercise}{[PX-C04.S03-06] Resuelve \((x-3)^3\geq 0\).}
\end{exercise}

\begin{shortanswer}{C04.S03 practica autonoma}
\begin{enumerate}[label=\textbf{PX-C04.S03-\arabic*:}, leftmargin=*]
  \item \((-\infty,-1)\cup (0,1)\)
  \item \([-3,+\infty)\)
  \item \((-\infty,0]\cup \{4\}\)
  \item \((-\infty,-2)\cup (5,+\infty)\)
  \item \((-\infty,1]\)
  \item \([3,+\infty)\)
\end{enumerate}
\end{shortanswer}

\begin{fullsolution}{C04.S03 practica autonoma}
\begin{enumerate}[label=\textbf{PX-C04.S03-\arabic*:}, leftmargin=*]
  \item El signo alterna en \(-1\), \(0\) y \(1\), quedando negativo en \((-\infty,-1)\) y \((0,1)\).
  \item \((x-2)^2\geq 0\), asi que el signo lo decide \(x+3\); se incluye \(-3\).
  \item El signo lo decide \(x\), pero en \(x=4\) el producto vale \(0\).
  \item Producto positivo fuera de las raices.
  \item Como \(x^2\geq 0\), el signo depende de \(x-1\) y el producto es no positivo si \(x\leq 1\).
  \item Una potencia impar conserva el signo de \(x-3\).
\end{enumerate}
\end{fullsolution}

\section{[C04.S04] Inecuaciones racionales}

\subsection*{Objetivos de aprendizaje}
\begin{itemize}
  \item Resolver inecuaciones con cocientes de polinomios.
  \item Separar raices del numerador y puntos prohibidos del denominador.
\end{itemize}

\subsection*{Prerrequisitos}
\begin{itemize}
  \item Inecuaciones polinomicas y fracciones algebraicas.
\end{itemize}

\begin{theorybox}{Idea minima}
En una inecuacion racional, los puntos importantes son las raices del numerador y los valores
que anulan el denominador. Estos ultimos nunca pueden pertenecer a la solucion.
\end{theorybox}

\begin{methodbox}{Como reconocer y resolver}
\begin{enumerate}
  \item Factoriza numerador y denominador.
  \item Marca las raices y los puntos prohibidos.
  \item Estudia el signo en cada intervalo.
  \item Incluye solo los ceros del numerador cuando la desigualdad lo permita.
\end{enumerate}
\end{methodbox}

\begin{solvedexample}{[EX-C04.S04-01] Cociente con un punto prohibido}
Resuelve
\[
\frac{x-3}{x+1}\leq 0.
\]

\textbf{Desarrollo.} El numerador se anula en \(x=3\) y el denominador en \(x=-1\), que es
punto prohibido. El estudio de signos da:
\[
(-\infty,-1): +,\qquad (-1,3): -,\qquad (3,+\infty): +.
\]
Como queremos \(\leq 0\), tomamos el intervalo negativo y el cero del numerador.

\textbf{Respuesta final.}
\[
(-1,3].
\]
\end{solvedexample}

\begin{commonerror}{Incluir el punto prohibido}
El valor que anula el denominador nunca se incluye, aunque la desigualdad sea \(\leq\) o \(\geq\).
\end{commonerror}

\begin{guidedexercise}{[GX-C04.S04-01] Signo positivo}
Resuelve
\[
\frac{x}{x-2}>0.
\]
\end{guidedexercise}

\begin{shortanswer}{[GX-C04.S04-01]}
\[
(-\infty,0)\cup (2,+\infty).
\]
\end{shortanswer}

\begin{fullsolution}{[GX-C04.S04-01]}
El cociente es positivo cuando numerador y denominador tienen el mismo signo:
\[
(-\infty,0)\cup (2,+\infty).
\]
\end{fullsolution}

\begin{guidedexercise}{[GX-C04.S04-02] Signo no negativo}
Resuelve
\[
\frac{x+1}{x}\geq 0.
\]
\end{guidedexercise}

\begin{shortanswer}{[GX-C04.S04-02]}
\[
(-\infty,-1]\cup (0,+\infty).
\]
\end{shortanswer}

\begin{fullsolution}{[GX-C04.S04-02]}
Los puntos clave son \(-1\) y \(0\). El cociente es no negativo en
\[
(-\infty,-1]\cup (0,+\infty).
\]
\end{fullsolution}

\subsection*{Practica autonoma graduada}
\begin{exercise}{[PX-C04.S04-01] Resuelve \(\dfrac{1}{x-4}<0\).}
\end{exercise}
\begin{exercise}{[PX-C04.S04-02] Resuelve \(\dfrac{x-2}{x+2}\geq 0\).}
\end{exercise}
\begin{exercise}{[PX-C04.S04-03] Resuelve \(\dfrac{x}{x+1}\leq 0\).}
\end{exercise}
\begin{exercise}{[PX-C04.S04-04] Resuelve \(\dfrac{x+3}{x-1}>0\).}
\end{exercise}
\begin{exercise}{[PX-C04.S04-05] Resuelve \(\dfrac{x-1}{x-3}<0\).}
\end{exercise}
\begin{exercise}{[PX-C04.S04-06] Resuelve \(\dfrac{1}{x^2-4}>0\).}
\end{exercise}

\begin{shortanswer}{C04.S04 practica autonoma}
\begin{enumerate}[label=\textbf{PX-C04.S04-\arabic*:}, leftmargin=*]
  \item \((-\infty,4)\)
  \item \((-\infty,-2)\cup [2,+\infty)\)
  \item \((-1,0]\)
  \item \((-\infty,-3)\cup (1,+\infty)\)
  \item \((1,3)\)
  \item \((-\infty,-2)\cup (2,+\infty)\)
\end{enumerate}
\end{shortanswer}

\begin{fullsolution}{C04.S04 practica autonoma}
\begin{enumerate}[label=\textbf{PX-C04.S04-\arabic*:}, leftmargin=*]
  \item El signo lo decide \(x-4\); la fraccion es negativa si \(x<4\).
  \item Mismo signo en numerador y denominador, incluyendo \(x=2\).
  \item Entre \(-1\) y \(0\) el cociente es no positivo, e incluye \(x=0\).
  \item Positiva si ambos factores tienen el mismo signo, sin incluir \(x=-3\) porque hace cero el numerador.
  \item Negativa entre las raices y sin incluirlas.
  \item \(x^2-4=(x-2)(x+2)\); el producto es positivo fuera de \([-2,2]\), sin incluir los extremos.
\end{enumerate}
\end{fullsolution}

\section{[C04.S05] Sistemas de inecuaciones en una variable}

\subsection*{Objetivos de aprendizaje}
\begin{itemize}
  \item Resolver varias inecuaciones a la vez.
  \item Hallar la interseccion final de sus soluciones.
\end{itemize}

\subsection*{Prerrequisitos}
\begin{itemize}
  \item Inecuaciones lineales y expresion en intervalos.
\end{itemize}

\begin{theorybox}{Idea minima}
Un sistema de inecuaciones se resuelve encontrando la solucion de cada una por separado y
tomando la interseccion de todas ellas.
\end{theorybox}

\begin{methodbox}{Como reconocer y resolver}
\begin{enumerate}
  \item Resuelve cada inecuacion por separado.
  \item Escribe cada solucion como intervalo.
  \item Intersecta los intervalos.
\end{enumerate}
\end{methodbox}

\begin{solvedexample}{[EX-C04.S05-01] Interseccion de dos condiciones}
Resuelve
\[
\begin{cases}
x+2>0,\\
2x-3\leq 1.
\end{cases}
\]

\textbf{Desarrollo.}
\[
x>-2,
\qquad
2x\leq 4 \Longrightarrow x\leq 2.
\]
La interseccion es
\[
(-2,2].
\]

\textbf{Respuesta final.}
\[
x\in (-2,2].
\]
\end{solvedexample}

\begin{commonerror}{Unir en lugar de intersectar}
En un sistema deben cumplirse todas las condiciones, asi que se intersectan soluciones; no se
unen.
\end{commonerror}

\begin{guidedexercise}{[GX-C04.S05-01] Dos cotas}
Resuelve
\[
\begin{cases}
x<5,\\
x\geq 1.
\end{cases}
\]
\end{guidedexercise}

\begin{shortanswer}{[GX-C04.S05-01]}
\[
[1,5).
\]
\end{shortanswer}

\begin{fullsolution}{[GX-C04.S05-01]}
La primera condicion da \((-\infty,5)\) y la segunda \([1,+\infty)\). Su interseccion es
\[
[1,5).
\]
\end{fullsolution}

\begin{guidedexercise}{[GX-C04.S05-02] Dos lineales}
Resuelve
\[
\begin{cases}
2x+1>3,\\
\dfrac{x}{2}\leq 4.
\end{cases}
\]
\end{guidedexercise}

\begin{shortanswer}{[GX-C04.S05-02]}
\[
(1,8].
\]
\end{shortanswer}

\begin{fullsolution}{[GX-C04.S05-02]}
\[
2x>2 \Longrightarrow x>1,
\qquad
x\leq 8.
\]
Por tanto,
\[
(1,8].
\]
\end{fullsolution}

\subsection*{Practica autonoma graduada}
\begin{exercise}{[PX-C04.S05-01] Resuelve \(\begin{cases}x>-3\\ x<4\end{cases}\).}
\end{exercise}
\begin{exercise}{[PX-C04.S05-02] Resuelve \(\begin{cases}x\geq 0\\ x\leq 6\end{cases}\).}
\end{exercise}
\begin{exercise}{[PX-C04.S05-03] Resuelve \(\begin{cases}x-1>0\\ 3-x\geq 0\end{cases}\).}
\end{exercise}
\begin{exercise}{[PX-C04.S05-04] Resuelve \(\begin{cases}2x\leq 8\\ x+5>0\end{cases}\).}
\end{exercise}
\begin{exercise}{[PX-C04.S05-05] Resuelve \(\begin{cases}x<2\\ x\geq 2\end{cases}\).}
\end{exercise}
\begin{exercise}{[PX-C04.S05-06] Resuelve \(\begin{cases}x\leq -1\\ x>-4\end{cases}\).}
\end{exercise}

\begin{shortanswer}{C04.S05 practica autonoma}
\begin{enumerate}[label=\textbf{PX-C04.S05-\arabic*:}, leftmargin=*]
  \item \((-3,4)\)
  \item \([0,6]\)
  \item \((1,3]\)
  \item \((-5,4]\)
  \item \(\varnothing\)
  \item \((-4,-1]\)
\end{enumerate}
\end{shortanswer}

\begin{fullsolution}{C04.S05 practica autonoma}
\begin{enumerate}[label=\textbf{PX-C04.S05-\arabic*:}, leftmargin=*]
  \item Interseccion de \(( -3,+\infty)\) y \((-\infty,4)\).
  \item Interseccion de \([0,+\infty)\) y \((-\infty,6]\).
  \item \(x>1\) y \(x\leq 3\).
  \item \(x\leq 4\) y \(x>-5\).
  \item No existe ningun \(x\) que cumpla ambas condiciones.
  \item \(x\in(-4,-1]\).
\end{enumerate}
\end{fullsolution}

\section{[C04.S06] Inecuaciones de dos variables y semiplanos}

\subsection*{Objetivos de aprendizaje}
\begin{itemize}
  \item Identificar la recta frontera de una inecuacion lineal en dos variables.
  \item Decidir que semiplano satisface la condicion.
\end{itemize}

\subsection*{Prerrequisitos}
\begin{itemize}
  \item Ecuacion de la recta y representacion en el plano.
\end{itemize}

\begin{theorybox}{Idea minima}
La inecuacion \(ax+by\leq c\) se apoya en la recta frontera \(ax+by=c\). Un punto de prueba
como \((0,0)\) suele ayudar a decidir que lado de la recta pertenece a la solucion.
\end{theorybox}

\begin{methodbox}{Como reconocer y resolver}
\begin{enumerate}
  \item Sustituye el signo por igualdad para obtener la recta frontera.
  \item Elige un punto sencillo que no este en la recta.
  \item Comprueba si ese punto satisface la inecuacion.
  \item Decide el semiplano correspondiente.
\end{enumerate}
\end{methodbox}

\begin{solvedexample}{[EX-C04.S06-01] Semiplano por punto de prueba}
Describe geometricamente la solucion de
\[
x+2y\leq 4.
\]

\textbf{Desarrollo.} La recta frontera es
\[
x+2y=4.
\]
Probamos el punto \((0,0)\):
\[
0+2\cdot 0\leq 4,
\]
que es verdadero. Por tanto, la solucion es el semiplano que contiene al origen, con la recta
incluida.

\textbf{Respuesta final.}
Recta frontera \(x+2y=4\) y semiplano que contiene \((0,0)\).
\end{solvedexample}

\begin{commonerror}{Cambiar el lado equivocado}
No basta con dibujar la recta: hay que comprobar con un punto de prueba que lado cumple la
inecuacion.
\end{commonerror}

\begin{guidedexercise}{[GX-C04.S06-01] Semiplano superior}
Describe la solucion de
\[
y>x.
\]
\end{guidedexercise}

\begin{shortanswer}{[GX-C04.S06-01]}
Recta frontera \(y=x\); semiplano superior, que contiene por ejemplo \((0,1)\).
\end{shortanswer}

\begin{fullsolution}{[GX-C04.S06-01]}
La frontera es \(y=x\). El punto \((0,1)\) cumple \(1>0\), asi que la solucion es el semiplano
por encima de la recta.
\end{fullsolution}

\begin{guidedexercise}{[GX-C04.S06-02] Frontera vertical}
Describe la solucion de
\[
x\geq -1.
\]
\end{guidedexercise}

\begin{shortanswer}{[GX-C04.S06-02]}
Recta frontera \(x=-1\); semiplano a la derecha, que incluye por ejemplo \((0,0)\).
\end{shortanswer}

\begin{fullsolution}{[GX-C04.S06-02]}
La recta frontera es vertical:
\[
x=-1.
\]
Como el punto \((0,0)\) verifica \(0\geq -1\), la solucion es el semiplano derecho, incluyendo
la recta.
\end{fullsolution}

\subsection*{Practica autonoma graduada}
\begin{exercise}{[PX-C04.S06-01] Describe la solucion de \(x+y\leq 3\).}
\end{exercise}
\begin{exercise}{[PX-C04.S06-02] Describe la solucion de \(y\geq 2\).}
\end{exercise}
\begin{exercise}{[PX-C04.S06-03] Describe la solucion de \(x<4\).}
\end{exercise}
\begin{exercise}{[PX-C04.S06-04] Describe la solucion de \(2x-y>0\).}
\end{exercise}
\begin{exercise}{[PX-C04.S06-05] Describe la solucion de \(x+y\geq 0\).}
\end{exercise}
\begin{exercise}{[PX-C04.S06-06] Describe la solucion de \(y\leq -x+2\).}
\end{exercise}

\begin{shortanswer}{C04.S06 practica autonoma}
\begin{enumerate}[label=\textbf{PX-C04.S06-\arabic*:}, leftmargin=*]
  \item Recta \(x+y=3\); semiplano que contiene \((0,0)\)
  \item Recta \(y=2\); semiplano superior, que contiene \((0,3)\)
  \item Recta \(x=4\); semiplano izquierdo, que contiene \((0,0)\)
  \item Recta \(y=2x\); semiplano inferior, que contiene \((1,0)\)
  \item Recta \(x+y=0\); semiplano que contiene \((1,0)\)
  \item Recta \(y=-x+2\); semiplano inferior, que contiene \((0,0)\)
\end{enumerate}
\end{shortanswer}

\begin{fullsolution}{C04.S06 practica autonoma}
\begin{enumerate}[label=\textbf{PX-C04.S06-\arabic*:}, leftmargin=*]
  \item La frontera es \(x+y=3\). El origen satisface \(0\leq 3\).
  \item La frontera es horizontal: \(y=2\). Se toman los puntos con ordenada mayor o igual que \(2\).
  \item La frontera es \(x=4\). Como \(0<4\), se toma el lado izquierdo.
  \item \(2x-y>0\) equivale a \(y<2x\); la frontera es \(y=2x\).
  \item La frontera es \(x+y=0\). El punto \((1,0)\) cumple \(1\geq 0\).
  \item Ya esta despejada: \(y\leq -x+2\), es el semiplano por debajo de la recta.
\end{enumerate}
\end{fullsolution}

\section{[C04.S07] Sistemas de inecuaciones en el plano}

\subsection*{Objetivos de aprendizaje}
\begin{itemize}
  \item Combinar varias inecuaciones lineales en el plano.
  \item Describir la region solucion mediante sus vertices o su forma.
\end{itemize}

\subsection*{Prerrequisitos}
\begin{itemize}
  \item Semiplanos y ecuacion de la recta.
\end{itemize}

\begin{theorybox}{Idea minima}
La region solucion de un sistema de inecuaciones en el plano es la interseccion de varios
semiplanos. Cuando la region es poligonal, sus vertices ayudan a describirla con precision.
\end{theorybox}

\begin{methodbox}{Como reconocer y resolver}
\begin{enumerate}
  \item Representa o identifica cada semiplano.
  \item Busca la zona comun.
  \item Si la region es acotada, localiza sus vertices.
\end{enumerate}
\end{methodbox}

\begin{solvedexample}{[EX-C04.S07-01] Triangulo en el primer cuadrante}
Describe la region solucion de
\[
\begin{cases}
x\geq 0,\\
y\geq 0,\\
x+y\leq 4.
\end{cases}
\]

\textbf{Desarrollo.} Las dos primeras condiciones nos dejan en el primer cuadrante. La tercera
anade el semiplano por debajo de la recta \(x+y=4\). La region comun es un triangulo.

\textbf{Respuesta final.}
Triangulo con vertices
\[
(0,0),\qquad (4,0),\qquad (0,4).
\]
\end{solvedexample}

\begin{commonerror}{Olvidar una condicion}
Si se dibuja solo una parte del sistema, la region obtenida puede ser mucho mas grande que la
solucion real.
\end{commonerror}

\begin{guidedexercise}{[GX-C04.S07-01] Triangulo desplazado}
Describe la region solucion de
\[
\begin{cases}
x\geq 1,\\
y\geq 0,\\
x+y\leq 5.
\end{cases}
\]
\end{guidedexercise}

\begin{shortanswer}{[GX-C04.S07-01]}
Triangulo con vertices \((1,0)\), \((5,0)\) y \((1,4)\).
\end{shortanswer}

\begin{fullsolution}{[GX-C04.S07-01]}
La recta vertical \(x=1\), el eje \(x\) y la recta \(x+y=5\) delimitan un triangulo con esos
vertices.
\end{fullsolution}

\begin{guidedexercise}{[GX-C04.S07-02] Triangulo apoyado sobre \(y=1\)}
Describe la region solucion de
\[
\begin{cases}
x\geq 0,\\
y\geq 1,\\
x+y\leq 3.
\end{cases}
\]
\end{guidedexercise}

\begin{shortanswer}{[GX-C04.S07-02]}
Triangulo con vertices \((0,1)\), \((2,1)\) y \((0,3)\).
\end{shortanswer}

\begin{fullsolution}{[GX-C04.S07-02]}
La region comun queda entre la recta \(y=1\), el eje \(y\) y la recta \(x+y=3\). Los cortes dan
los vertices indicados.
\end{fullsolution}

\subsection*{Practica autonoma graduada}
\begin{exercise}{[PX-C04.S07-01] Describe la region de \(\begin{cases}x\geq 0\\ y\geq 0\\ x+y\leq 2\end{cases}\).}
\end{exercise}
\begin{exercise}{[PX-C04.S07-02] Describe la region de \(\begin{cases}x\geq -1\\ y\geq 0\\ x+y\leq 3\end{cases}\).}
\end{exercise}
\begin{exercise}{[PX-C04.S07-03] Describe la region de \(\begin{cases}x\geq 0\\ y\geq 0\\ x\leq 3\\ y\leq 2\end{cases}\).}
\end{exercise}
\begin{exercise}{[PX-C04.S07-04] Describe la region de \(\begin{cases}x\geq 0\\ y\geq 0\\ y\leq x\\ x+y\leq 4\end{cases}\).}
\end{exercise}
\begin{exercise}{[PX-C04.S07-05] Describe la region de \(\begin{cases}x\geq 1\\ y\geq 1\\ x+y\leq 5\end{cases}\).}
\end{exercise}
\begin{exercise}{[PX-C04.S07-06] Describe la region de \(\begin{cases}x\geq 0\\ y\geq 0\\ x+2y\leq 4\end{cases}\).}
\end{exercise}

\begin{shortanswer}{C04.S07 practica autonoma}
\begin{enumerate}[label=\textbf{PX-C04.S07-\arabic*:}, leftmargin=*]
  \item Triangulo con vertices \((0,0)\), \((2,0)\), \((0,2)\)
  \item Triangulo con vertices \((-1,0)\), \((3,0)\), \((-1,4)\)
  \item Rectangulo con vertices \((0,0)\), \((3,0)\), \((3,2)\), \((0,2)\)
  \item Triangulo con vertices \((0,0)\), \((4,0)\), \((2,2)\)
  \item Triangulo con vertices \((1,1)\), \((4,1)\), \((1,4)\)
  \item Triangulo con vertices \((0,0)\), \((4,0)\), \((0,2)\)
\end{enumerate}
\end{shortanswer}

\begin{fullsolution}{C04.S07 practica autonoma}
\begin{enumerate}[label=\textbf{PX-C04.S07-\arabic*:}, leftmargin=*]
  \item Primer cuadrante por debajo de \(x+y=2\).
  \item La recta \(x=-1\), el eje \(x\) y \(x+y=3\) recortan ese triangulo.
  \item Interseccion de cuatro semiplanos: banda horizontal y banda vertical.
  \item La recta \(y=x\) corta a \(x+y=4\) en \((2,2)\), y el eje \(x\) aporta \((0,0)\) y \((4,0)\).
  \item Desplazando el primer cuadrante al punto \((1,1)\) y cortando con \(x+y=5\).
  \item Interseccion con los ejes y la recta \(x+2y=4\).
\end{enumerate}
\end{fullsolution}

\begin{selfcheck}
\begin{itemize}
  \item Distingo bien entre union e interseccion cuando trabajo con varias condiciones.
  \item En inecuaciones racionales no olvido los puntos prohibidos.
  \item En el plano, describo la frontera y luego justifico el lado correcto.
\end{itemize}
\end{selfcheck}

\begin{extension}{Leer la geometria de una desigualdad}
Una inecuacion en el plano no es solo un calculo: representa una region. Saber leer esa region
ayuda mucho despues en optimizacion y programacion lineal.
\end{extension}
